\subsection*{The distribution of represented orientations}
Suppose \(n\ge2\) is even and \(L\) is positive definite. Let \(U(m)\) and \(T(m)\) count chiral ordinary classes represented with exactly one and exactly two proper orientations, respectively. Then
\[
\begin{aligned}
 \sum_{m\le X}U(m)&=\frac{2R_O-R_{SO}}{n}X^n+
                         O_{L,\epsilon}(X^{n-1+\epsilon}),\\
 \sum_{m\le X}T(m)&=\frac{R_{SO}-R_O}{n}X^n+
                         O_{L,\epsilon}(X^{n-1+\epsilon}).
\end{aligned}
\]
The two coefficient identities are \(a_O=b^-+U+T\) and
\(a_{SO}=b^-+U+2T\). In every even dimension the improper transporter coefficient is \(b^-/2\), including in rank two, so the absolute convergence of the nonscalar remainder in Corollary~\ref{cor:KLL} bounds its summatory function by \(O_{L,\epsilon}(X^{n-1+\epsilon})\). Apply Theorem~\ref{th:analytic}. In particular a chiral parent represents a positive limiting proportion of each kind among all ordinary classes ordered by index. For an achiral parent \(U(m)=0\) for every \(m\).

For either orientation group in any positive rank \(n\ge2\), all but \(O_{L,\epsilon}(X^{n-1+\epsilon})\) represented classes of index at most \(X\) have \(\Gamma_H(M)=H\cap\{1,-1\}\). Indeed each exceptional class contributes \(1-c_H/|\Gamma_H(M)|\ge1/2\) to the nonscalar transporter sum. For ordinary equivalence such classes therefore have zero limiting proportion.

